%======================================================================
% CAPÍTULO 06: MÓDULO DE CONFIGURACIÓN
%======================================================================
% Parámetros del sistema, tasas BCV y directivas
%======================================================================

\chapter{Módulo de Configuración}
\label{ch:configuracion}

\section{Descripción General}

El Módulo de Configuración permite la administración de los parámetros generales del sistema SSSIFANB. Este módulo está restringido a usuarios con rol de administrador y contiene las opciones necesarias para mantener el correcto funcionamiento de la plataforma.

\section{Submódulos Disponibles}

\begin{itemize}[leftmargin=*]
    \item \textbf{Parámetros}: Variables globales del sistema
    \item \textbf{Tasa BCV}: Cotización del dólar oficial
    \item \textbf{Directivas}: Gestión de documentos normativos
    \item \textbf{Auth Data}: Configuración de seguridad
\end{itemize}

\section{Parámetros del Sistema}

\subsection*{Variables Globales}

\begin{longtable}{p{4cm}p{8cm}}
\toprule
\textbf{Parámetro} & \textbf{Descripción} \\
\midrule
\endhead
\bottomrule
\caption{Parámetros del sistema}
\label{tab:parametros}
\endlastfoot

\textbf{Tasa de Aporte} & Porcentaje de cotización mensual obligatoria \\
\textbf{Tope de Anticipo} & Límite máximo de anticipo permitted \\
\textbf{Porcentaje Fondo Social} & Asignación para programas sociales \\
\textbf{Interés Fideicomiso} & Tasa de rendimiento anual del fideicomiso \\
\textbf{Fecha de Cierre} & Día de corte para procesamiento de nómina \\
\end{longtable}

\subsection*{Cron Operativo Transversal (Homologación - Art. 39)}

Es vital para el sistema codificar que las actualizaciones ejecutadas sobre Parámetros detonen transaccionalmente un proceso de actualización sobre los expedientes del módulo \emph{Pensionados y Retroactivos}, dado que legalmente, toda pensión es equivalente al alza real activa de forma inmediata sin esperar comprobación extra.

\begin{lstlisting}[language=Java, caption=Trigger de Homologación Salarial]
function actualizarSueldoBaseDecreto(nuevoSueldoGeneral) {
    // Al guardarse el parametro de aumento, el CRON actualiza los expedientes 
    // sin requerir validaciones burocraticas manuales (Art. 39).
    triggerActualizacionMasiva({
        moduloDestino: 'Pensionados_Retroactivos',
        action: 'EQUIPARAR_ESCALA_EXPRESS',
        newBase: nuevoSueldoGeneral
    });
}
\end{lstlisting}

\subsection*{Edición de Parámetros}

Para modificar un parámetro:

\begin{enumerate}[leftmargin=*]
    \item Seleccionar \textbf{Parámetros} en el menú
    \item Buscar el parámetro a modificar
    \item Ingresar el nuevo valor
    \item Justificar el cambio en el campo de observaciones
    \item Hacer clic en \textbf{Guardar Cambios}
    \item Confirmar la acción en el diálogo de seguridad
\end{enumerate}

\begin{figure}[H]
\centering
\begin{tcolorbox}[
    enhanced,
    colback=sandraSoft,
    colframe=sandraAccent,
    width=0.9\textwidth,
    halign=center
]
\small
\textbf{Advertencia}\\
\smallskip
Los cambios en los parámetros del sistema afectan directamente los cálculos de nómina y prestaciones. Asegúrese de verificar la información antes de guardar y documente cualquier modificación para auditoría.
\end{tcolorbox}
\end{figure}

\section{Tasa del Banco Central de Venezuela}

\subsection*{Gestión de Tasas}

El sistema permite consultar y registrar la tasa oficial del BCV para conversiones monetarias:

\begin{figure}[H]
\centering
\begin{tcolorbox}[
    enhanced,
    colback=white,
    colframe=sandraDark,
    width=0.85\textwidth,
    halign=center
]
\begin{tabular}{lr}
\toprule
\textbf{Parámetro} & \textbf{Valor} \\
\midrule
Moneda Base & Bolívar (VES) \\
Moneda Extranjera & Dólar (USD) \\
Tasa BCV Actual & \$\$85,25 por USD \\
Fecha de Actualización & 15/03/2026 \\
Hora de Actualización & 10:00 AM \\
\midrule
\textbf{Última Compra} & \$\$85,18 \\
\textbf{Última Venta} & \$\$85,32 \\
\bottomrule
\end{tabular}
\end{tcolorbox}
\caption{Tasa BCV actual}
\label{fig:tasa-bcv}
\end{figure}

\subsection*{Sincronización Automática}

\begin{itemize}[leftmargin=*]
    \item El sistema puede configurarse para sincronizar automáticamente la tasa
    \item La actualización se realiza mediante conexión con el API del BCV
    \item Se mantiene registro histórico de todas las tasas aplicadas
    \item Las conversiones usan la tasa del día de procesamiento
\end{itemize}

\section{Directivas PDF}

\subsection*{Gestión de Documentos}

Este submódulo permite cargar y administrar los documentos normativos del sistema:

\begin{figure}[H]
\begin{tcolorbox}[
    enhanced,
    colback=sandraSoft,
    colframe=sandraDark,
    width=0.9\textwidth,
    halign=center
]
\small
\textbf{Tipos de Directivas}\\
\smallskip
\begin{itemize}
    \item Leyes y reglamentos
    \item Circulares ministeriales
    \item Resoluciones internas
    \item Instructivos operativos
    \item Formatos oficiales
\end{itemize}
\end{tcolorbox}
\end{figure}

\subsection*{Cargar Nueva Directriz}

\begin{enumerate}[leftmargin=*]
    \item Seleccionar \textbf{Directivas} en el menú
    \item Hacer clic en \textbf{Agregar Directriz}
    \item Completar los metadatos:
    \begin{itemize}
        \item Número de documento
        \item Título
        \item Fecha de vigencia
        \item Categoría
        \item Descripción
    \end{itemize}
    \item Adjuntar archivo PDF
    \item Hacer clic en \textbf{Guardar}
\end{enumerate}

\subsection*{Especificaciones de Archivos}

\begin{longtable}{p{3cm}p{9cm}}
\toprule
\textbf{Parámetro} & \textbf{Especificación} \\
\midrule
\endhead
\bottomrule
\caption{Especificaciones para carga de directivas}
\label{tab:especificaciones-directivas}
\endlastfoot

\textbf{Formato} & PDF exclusivamente \\
\textbf{Tamaño máximo} & 25 MB \\
\textbf{Resolución} & Mínimo 150 DPI para legibilidad \\
\textbf{Nomenclatura} & \texttt{TIPO\_NUMERO\_AÑO.pdf} \\
\end{longtable}

\section{Auth Data (Seguridad)}

\subsection*{Configuración de Autenticación}

Este submódulo gestiona los parámetros de seguridad del ecosistema:

\begin{itemize}[leftmargin=*]
    \item \textbf{Habilitar/Deshabilitar 2FA}: Control global de autenticación de dos factores
    \item \textbf{Método TOTP}: Configuración del tiempo de validez de códigos
    \item \textbf{Límite de intentos}: Número máximo de intentos de login fallidos
    \item \textbf{Tiempo de bloqueo}: Duración del bloqueo por intentos excesivos
    \item \textbf{Vencimiento de sesión}: Tiempo máximo de inactividad
\end{itemize}

\begin{figure}[H]
\centering
\begin{tcolorbox}[
    enhanced,
    colback=sandraSoft,
    colframe=sandraGreen,
    width=0.9\textwidth,
    halign=center
]
\small
\textbf{Seguridad del Sistema}\\
\smallskip
El SSSIFANB implementa múltiples capas de seguridad incluyendo encriptación HMAC SHA256 para contraseñas, tokens JWT con expiración configurable, y autenticación de dos factores mediante códigos TOTP. Los cambios en esta configuración deben ser realizados por personal autorizado.
\end{tcolorbox}
\end{figure}

\section{Logs de Auditoría}

El sistema mantiene un registro de todas las modificaciones realizadas:

\begin{longtable}{p{3cm}p{3cm}p{3cm}p{4cm}}
\toprule
\textbf{Fecha} & \textbf{Hora} & \textbf{Usuario} & \textbf{Acción} \\
\midrule
\endhead
\bottomrule
\caption{Ejemplo de log de auditoría}
\label{tab:log-auditoria}
\endlastfoot

15/03/2026 & 09:45 & admin\_ipsfa & Modificó Tasa de Aporte \\
15/03/2026 & 10:12 & admin\_ipsfa & Cargó Directiva N° 2854 \\
15/03/2026 & 11:30 & admin\_ipsfa & Habilitó 2FA global \\
16/03/2026 & 08:00 & admin\_ipsfa & Actualizó tasa BCV \\
\end{longtable}

\section{Bitácora de Cambios}

\begin{enumerate}[leftmargin=*]
    \item Seleccionar \textbf{Ver Bitácora} en cualquier submódulo
    \item Filtrar por período, usuario o tipo de acción
    \item Exportar reporte en PDF o Excel
    \item Consultar detalles de cada modificación
\end{enumerate}

\begin{figure}[H]
\begin{tcolorbox}[
    enhanced,
    colback=sandraSoft,
    colframe=sandraDark,
    width=0.9\textwidth,
    halign=center
]
\small
\textbf{Retención de Logs}\\
\smallskip
Los registros de auditoría se mantienen durante 5 años según las políticas de retención de datos del IPSFANB. Los logs pueden ser solicitados por organismos de control en cualquier momento.
\end{tcolorbox}
\end{figure}

\section{Respaldo y Restauración}

\subsection*{Respaldo de Base de Datos}

El administrador puede generar respaldos de la base de datos:

\begin{enumerate}[leftmargin=*]
    \item Acceder a \textbf{Configuración > Respaldo}
    \item Seleccionar el tipo de respaldo:
    \begin{itemize}
        \item Completo: Toda la base de datos
        \item Parcial: Solo tablas seleccionadas
    \end{itemize}
    \item Hacer clic en \textbf{Generar Respaldo}
    \item Descargar el archivo de respaldo generado
    \item Almacenar en ubicación segura
\end{enumerate}

\subsection*{Restauración}

\begin{enumerate}[leftmargin=*]
    \item Acceder a \textbf{Configuración > Restauración}
    \item Seleccionar el archivo de respaldo
    \item Verificar integridad del archivo
    \item Confirmar la restauración
    \item Esperar proceso de importación
    \item Validar datos restaurados
\end{enumerate}

\begin{figure}[H]
\centering
\begin{tcolorbox}[
    enhanced,
    colback=sandraSoft,
    colframe=sandraAccent,
    width=0.9\textwidth,
    halign=center
]
\small
\textbf{Advertencia de Restauración}\\
\smallskip
La restauración de datos reemplaza la información actual. Este proceso debe realizarse únicamente en casos de contingencia y requiere confirmación escrita del director del IPSFANB.
\end{tcolorbox}
\end{figure}

\section{Usuarios y Permisos}

\subsection*{Gestión de Usuarios}

\begin{longtable}{p{3cm}p{9cm}}
\toprule
\textbf{Acción} & \textbf{Descripción} \\
\midrule
\endhead
\bottomrule
\caption{Operaciones de gestión de usuarios}
\label{tab:gestion-usuarios}
\endlastfoot

\textbf{Crear Usuario} & Agregar nuevo usuario al sistema \\
\textbf{Modificar Usuario} & Cambiar datos o permisos \\
\textbf{Desactivar Usuario} & Suspender acceso sin eliminar \\
\textbf{Eliminar Usuario} & Remover usuario del sistema \\
\textbf{Resetear Clave} & Forzar cambio de contraseña \\
\end{longtable}

\subsection*{Roles Disponibles}

\begin{itemize}[leftmargin=*]
    \item \textbf{Administrador}: Acceso total al sistema
    \item \textbf{Operador}: Gestión de solicitudes
    \item \textbf{Consulta}: Solo lectura de información
    \item \textbf{Afiliado}: Acceso a datos propios
\end{itemize}

\newpage
